\section{局限与复现}
\label{sec:limitations}
\subsection{结果没有建立什么}
历史材料是观察性、回顾性且经过筛选的。按任务归组保留了重复观察之间的联系，却不能消除选择偏差，也不能确认不同任务彼此独立。共同的数学依赖、审核者和执行条件仍可能带来依赖。统计区间采用第~\ref{sec:statistics}~节和第~\ref{sec:process}~节说明的参考解释。

要求检查与历史定义的解释均有智能体参与，没有经过独立专家参照的裁定。定义见证确认指定形式谓词之间的关系，布尔对照穷尽一个只有四种输入的任务；二者都不能给出全部历史的题意评价准确率。

可执行检查使用记录的 Lean 与 Mathlib 环境，没有完整重建每份候选最初的环境，因此当前编译失败未必是原有缺陷。见证由记录的 Lean 安装检查，没有采用第二套外部内核实现。部分重建步骤复用了原代码。独立统计复算检查冻结输入和公式，哈希值检查文件与所给清单是否一致；二者都不能确认外部真实性或随机抽样。

\subsection{本版范围}
第 0.5.0 版按实际工作及后续分析操作重组全文，恢复第 0.3.0 版的定义入口，保留第 0.4.0 版的平实行文，默认读者掌握常用数学记号和统计基础。既有论证、比较例、实证表格及推断边界完整保留；没有新增实验或拟合模型，此前版本未改。

\subsection{源文件与复现}
运行 \code{python -B build\_reports.py}，即可通过 Tectonic 编译 \artifactpath{latex/zh/} 与 \artifactpath{latex/en/} 中的 LaTeX 源稿及 \artifactpath{latex/figures/} 中的图件，输出两份 PDF 至 \artifactpath{output/pdf/}。字体与构建要求见 \artifactpath{README.md}；版式检查及目录迁移构建均使用所记录的 Windows 字体环境。

\artifactpath{analysis/} 保留第 0.2.0 版的冻结输入、脚本、结果和依赖要求；编译报告无需重跑分析。便携输入复现所规定的投影，原始全历史收集仍需原档案。分类与连续性来源、目录映射及公开仓库见附录~\ref{sec:sources}。交付记录将源码与 PDF 绑定到内容和版式检查。

\subsection{智能体辅助}
智能体参与了实现、部分题意解释、数学与统计核查及写作。此前的详解版采用了多个撰写角色，本次行文修订由主任务完成。这些检查不等于独立的人类同行评审。报告与源稿包按所链接评估仓库的分发结构准备。
